% ===================================================================
% SKILL: ÇİFT SÜTUNLU HİZALI KUTULAR (SIDE-BY-SIDE BOXES)
% ===================================================================
% TANIM:
% Beamer'da iki sütunlu (columns) tasarımda kutuların üst hizalamasının
% kayması çok yaygın bir sorundur. Bu şablon, [t] (top align)
% parametresi ile kutuların kusursuz bir şekilde aynı hizada
% başlamasını sağlar.
% 
% MEKANİZMA:
% \begin{columns}[t] -> Sütun içeriğini en üstten hizalar.
% ===================================================================

\begin{frame}{Karşılaştırmalı Analiz (Hizalı Kutular)}
    
    % [t] parametresi CRITICAL öneme sahiptir.
    % Olmazsa içeriğe göre dikeyde ortalanır ve hizalar kayar.
    \begin{columns}[t] 
        
        % SOL SÜTUN
        \column{0.48\textwidth}
        \begin{block}{1. Tanımlayıcı İstatistik}
            % İçerik az olsa bile kutu en üstten başlar.
            Verilerin özetlenmesi ve betimlenmesidir.
            \begin{itemize}
                \item Ortalama
                \item Medyan
                \item Mod
            \end{itemize}
        \end{block}

        % SAĞ SÜTUN
        \column{0.48\textwidth}
        \begin{block}{2. Çıkarımsal İstatistik}
            % İçerik çok olsa bile sol kutuyla aynı hizada başlar.
            Örneklem verisinden anakütle hakkında tahminlerde bulunmaktır.
            \begin{itemize}
                \item Hipotez Testleri
                \item Güven Aralıkları
                \item Regresyon Analizi
                \item ANOVA
            \end{itemize}
        \end{block}
        
    \end{columns}
    
    \vspace{1em}
    
    % FARKLI TÜR KUTULARIN YAN YANA KULLANIMI
    \begin{columns}[t]
        
        \column{0.48\textwidth}
        \begin{alertblock}{Yanlış Yöntem}
            Verinin sadece ortalamasına bakarak dağılım hakkında karar vermek yanıltıcıdır.
        \end{alertblock}
        
        \column{0.48\textwidth}
        \begin{exampleblock}{Doğru Yöntem}
            Histogram ve Kutu Grafiği (Boxplot) çizerek aykırı değerleri incelemektir.
        \end{exampleblock}
        
    \end{columns}

\end{frame}
